\chapter*[Introdução]{Introdução}
\addcontentsline{toc}{chapter}{Introdução}
Sistemas de machine learning são construídos sobre a premissa de que os dados
de produção se comportam como os dados de treinamento. O ambiente ao redor de
um modelo, porém, muda continuamente: usuários evoluem, mercados se
transformam, sensores se desgastam. Quando isso ocorre, um modelo bem
treinado passa a produzir resultados progressivamente piores sem emitir
qualquer erro explícito, fenômeno chamado de data drift, cuja característica
mais grave é o silêncio da degradação \cite{sculley2015}. \citeonline{gama2014}
formalizam o fenômeno como a mudança, ao longo do tempo, na distribuição
conjunta entre os dados de entrada e a variável alvo, comprometendo a validade
de um modelo treinado sobre uma distribuição que já não representa o presente.

Detectar quando essa mudança ocorreu é hoje um problema relativamente maduro,
com métodos que monitoram a distribuição das entradas ou o desempenho do
classificador \cite{bayram2022}. O que permanece pouco resolvido é a etapa
seguinte, a de decidir como reagir ao drift detectado. Uma prática comum é
adotar uma resposta única e fixa, tipicamente descartar o modelo e
retreiná-lo do zero. Há evidências, no entanto, de que nenhuma resposta é
ótima em todos os cenários: estratégias adequadas para real drift são
subótimas sob virtual drift \cite{oliveira2023}, o reaproveitamento de
modelos supera o retreinamento do zero quando conceitos reaparecem
\cite{sun2018}, e classificadores distintos reagem de formas substancialmente
diferentes ao mesmo drift \cite{webb2016}. Além do tipo de drift, sua
severidade também condiciona a resposta adequada: \citeonline{wang2022elastic},
em estudo sobre boosting adaptativo à severidade do drift,
demonstram formalmente que a severidade do drift determina a taxa de
convergência necessária para readaptação do modelo, e \citeonline{wang2024quadcdd}
argumentam que caracterizar um drift de forma completa exige representar tipo
e severidade conjuntamente, não apenas um desses eixos isoladamente. A
escolha da estratégia de resposta deveria, portanto, depender do tipo e da
severidade do drift, além do classificador em uso, dependência conjunta que
raramente é investigada de forma sistemática.

Diante desse cenário, este trabalho é orientado pela seguinte questão de
pesquisa: como a eficácia de diferentes estratégias de retreinamento varia em
função do tipo e da severidade do drift enfrentado e do classificador
monitorado, em ambiente de aprendizado em stream? Nessa questão, o tipo
refere-se especificamente ao padrão temporal do drift (abrupto, gradual,
incremental ou recorrente), não à distinção entre real drift e virtual
drift discutida acima, tratada neste trabalho como dimensão conceitual à
parte (Seção~\ref{sec:sinais}). A questão é respondida por
um experimento fatorial controlado, no qual o tipo e a severidade do drift
são conhecidos por injeção em dados sintéticos e o detector é fixo em todas
as condições, de modo a isolar a estratégia de resposta como variável de
interesse. Como desdobramento exploratório, investiga-se ainda se os padrões
observados permitem esboçar uma orientação preliminar que relacione sinais
observáveis em produção, tais como a magnitude da divergência calculada pelo
detector, usada como proxy de severidade, a velocidade com que essa
divergência varia, indicativa do padrão temporal, e a semelhança entre a
distribuição atual e distribuições históricas, indício de recorrência, a
estratégias de retreinamento recomendadas, de modo a orientar decisões
futuras em pipelines de MLOps (Machine Learning Operations).

Este trabalho está organizado da seguinte forma. O Capítulo 1 apresenta o
referencial teórico, estabelecendo a definição formal de drift, suas
dimensões de caracterização, os métodos de detecção, a operacionalização da
severidade, as estratégias de retreinamento, os classificadores considerados
e os sinais observáveis que fundamentam o desdobramento exploratório deste
trabalho e o Capítulo 2 descreve a metodologia, detalhando a geração da stream sintética, a adaptação dos classificadores ao ambiente de stream, a operacionalização
das estratégias de retreinamento, o desenho experimental fatorial, as
métricas de avaliação, o protocolo de validação da questão derivada e o
cronograma de execução.